\section{Materials-discovery applications and demonstrative tests}
\label{sec:applications}
\subsection{Hybrid-functional topological materials databases}
TopoHSE-DB provides an application example of electronic-structure-based
topological classification at database scale
\cite{Mirhosseini2026TopoHSE}. Mirhosseini, Elcoro, Kn\"upfer and
K\"uhne combine consistently relaxed nonmagnetic crystal structures,
PBE and screened-hybrid HSE06 electronic structures, and topological
quantum chemistry. Comparing the resulting labels exposes the
sensitivity of a symmetry-based classification to the underlying
exchange--correlation approximation. A converged integer or symmetry
label is therefore not, by itself, evidence that the electronic
Hamiltonian has been described accurately.

The original database workflow uses VASP, VASP2Trace, and
CheckTopologicalMat. It is cited here as a materials-discovery
application and a source of future comparison cases, not as a
calculation performed with the new CP2K topology routines.
Connecting such a workflow to CP2K requires equivalent electronic
structure and symmetry conventions as well as the general little-group,
compatibility-relation, and EBR machinery discussed in
Section~\ref{sec:tqc}. The currently implemented inversion analysis
does not yet reproduce the database's full classification workflow.

\subsection{Generative--predictive discovery of candidate materials}
The submitted study by Brocai, K\"uhne, Kn\"upfer and Mirhosseini,
\emph{Generative--predictive models enable the discovery of topological
materials}, provides a complementary application
\cite{Brocai2026Discovery}. It couples topology-conditioned crystal
generation to a predictive filter using a shared pretrained
representation, followed by hybrid-functional electronic-structure
and TQC analysis. The submitted manuscript reports 1,033 successfully
calculated hybrid-functional band structures; 629 receive an unambiguous
classification and 523 are classified as topological. The corresponding
83.15\% rate applies to the classified subset, whereas the confirmed
yield across all 1,033 calculated cases is 50.6\%. The 404 unresolved
classifications must not be silently counted as either trivial or
topological.

These results, obtained with the study's VASP-based workflow, illustrate
the demand for reliable post-processing and explicit failure diagnostics.
They also motivate subsequent, distinct questions: Wilson spectra can
probe isolated band subspaces, Kubo calculations can address transport
at specified temperature and dissipation, and localizers can investigate
finite, defective, or disordered realizations of selected candidates.
These are prospective CP2K follow-up applications, not calculations
already completed for the reported candidate set. A TQC label alone
does not supply a conductivity or a localizer protection gap.

\subsection{Small calculations as tests and demonstrations}
The smaller calculations in this paper and the companion k-point
study serve a different purpose \cite{KpointCompanion2026}.
They provide transparent tests and demonstrations of individual
algorithmic steps, not a second materials-discovery survey.
Table~\ref{tab:demonstration-roles} states the role of each set.
\begin{table}[htbp]
\centering\small
\begin{tabular}{L{0.27\linewidth}L{0.59\linewidth}}
\toprule
Calculation set & Demonstration and test objective\\
\midrule
Si and Al, companion k-point study &
Quadrature dependence and equivalent full-grid, time-reversal,
K290, and SPGLIB representations. At fixed $4^3$ sampling,
64 eigensystems reduce to four with atomic symmetry.\\
Periodic neon and buckled stanene &
Trivial and nontrivial Wilson-loop examples, with $\nu=0$ and
$\nu=1$, respectively; comparison of native and external analyses.\\
Ne/SOC, Bi$_2$/SOC, and diamond/SOC &
Primitive-mesh versus supercell and reduced versus full-property-grid
checks of the metric, Hamiltonian, and current operators.\\
Analytic localizer models and small GPW/GAPW systems &
Known indices, gap closings, basis covariance, and dense/sparse
solver consistency; quadratic-state residual and subspace checks.\\
\bottomrule
\end{tabular}
\caption{Test and demonstration calculations, distinct from the
materials-discovery applications cited above. The Si/Al results are
retained from the companion study, not rerun or reassigned to the
new topology implementation.}
\label{tab:demonstration-roles}
\end{table}

For Si and Al, the largest fixed-mesh cell-energy or electronic-free-energy
difference among the compared representations is
$1.4\times10^{-12}$ hartree. This tests representation equivalence,
not complete basis or k-point convergence. In particular, the Al
free energy still changes by 10.721 meV per atom between the
$6^3$ and $8^3$ meshes. Likewise, agreement between Wilson analyses
using the same overlap matrices tests the analysis stage, while
independent operator identities test the AO construction.
The quantitative checks in Section~\ref{sec:validation} and the
Supporting Information retain these distinctions and their original
calculation provenance.
